\section{Related Work}
\label{sec:relwork}
Our work sits at the intersection of edge IDS, conditional inference, model serving, and learned resource orchestration. Recent IDS research covers continual detection, heterogeneous representations, and runtime model selection~\cite{Wang2024BFSNID,AOCIDS2024,StatAvg2025,DMIDS2025}. Auto-IDaS and FEDMS show that learned or confidence-driven detector selection is operationally plausible~\cite{Lin2024AutoIDaS,Hasan2026FEDMS}, but neither a complex controller nor a cheaper operating point by itself establishes adaptive value. We therefore require comparison with the static detector actually realized by the selector and expose the full state--action map, not only visited-state occupancy.

Conditional-computation systems---including BranchyNet, adaptive network selection, MSDNet, and shallow--deep networks---escalate uncertain inputs to costlier models~\cite{Teerapittayanon2016BranchyNet,Bolukbasi2017Adaptive,Huang2018MSDNet,Kaya2019ShallowDeep}. Our fixed confidence cascade translates that idea to whole-device IDS measurement. It also exposes a serving detail often hidden by row-level escalation rates: only 1.34\% of rows escalate, yet they activate MedRF in 66.76\% of windows, so call granularity rather than row fraction governs the observed cost.

Serving systems such as Clipper, INFaaS, and Clockwork make accuracy, latency, and cost SLOs explicit and route requests among model variants~\cite{Crankshaw2017Clipper,Romero2021INFaaS,Gujarati2020Clockwork}. We adopt their service-contract perspective but add two missing controls: a matched static action that separates routing value from model choice, and a pre-deployment audit that asks whether the declared reward permits any action reversal on validation support.

RL-based edge orchestration spans vRAN, vehicle--edge scheduling, DVFS, offloading, and QoS control~\cite{Murti2024vRAN,Zhang2024VehicleEdge,Zhang2024DVFO,Lou2024Startup,Kim2024DistributedOffload,Maleki2024QCS,Nan2026RobustDNN,Wang2025TrustOffload,Hu2025SITOff}. Tabular Q-learning and DQN are representative discrete and neural value-based controllers~\cite{Watkins1992,Mnih2015}. Our contribution is not a new RL algorithm. It is a diagnostic methodology that combines an analytic reward screen, matched physical attribution, a prospectively frozen load/prevalence factorial with an explicit deadline SLO, and off-support action auditing. These instruments address known evaluation fragility in security ML and small-sample RL studies~\cite{Arp2022,Agarwal2021}.
